\begin{abstract}
Agent memory systems often ask a language model to decide what to remember before future questions are known. This adds ingestion cost and can discard details that later matter. We investigate whether effective memory requires this generative write step. Our approach preserves the original history and retrieves evidence when a question arrives, with no generative model calls to construct the store: conversation turns, or whole agent trajectories, are kept as they were, indexed by keyword and by a local embedding model, and reordered at question time by a fused score, two dialogue rules and two local cross-encoders. An agent pays a memory's build and read costs on every session while a benchmark score is paid once, so we report accuracy together with tokens read, search time and build cost. With GPT-5 reading the top 40 stored turns, the store answers 90.6 percent of 607 LoCoMo questions, 385 of them development questions used to tune the search, and 82.0 percent of the 222 test questions completed, at 3,294 reader tokens per question and about one second of search. Retrieval holds up without model-written memories: the store finds every gold session in the top 10 for 97.4 percent of 470 LongMemEval-S questions, including development questions used for tuning, recovers 80.3 percent of LoCoMo gold turns in the top 10 against 64.0 for a flat retrieval baseline, and ranks agent trajectories on Proced-Mem at 0.957 mean average precision. The remaining errors expose a tradeoff between preserving information and finding all the evidence needed to answer. These results motivate agent memory that retains the evidence and leaves its interpretation to the moment of use. The code, the per-question results and the protocol that registered every run are public.
\end{abstract}
